\documentclass[a4paper,11pt]{article}
\usepackage[T2A]{fontenc}
\usepackage[utf8]{inputenc}
\usepackage[russian]{babel}
\usepackage{amsmath,amssymb,mathtools}
\usepackage{geometry}
\geometry{a4paper,left=21mm,right=21mm,top=21mm,bottom=21mm}
\usepackage{microtype}
\usepackage{needspace}
\usepackage{tikz}
\usetikzlibrary{calc,arrows.meta,positioning,shapes.geometric}
\usepackage{booktabs,tabularx}
\usepackage{enumitem}
\usepackage{fancyhdr}
\usepackage{setspace}
\usepackage{xcolor}
\usepackage{hyperref}
\hypersetup{hidelinks,pdftitle={Чек-лист. Квадратные уравнения и алгебраические преобразования. Путь Б},pdfauthor={Лёвин Артём Александрович}}
\definecolor{accent}{HTML}{087C83}
\definecolor{darkaccent}{HTML}{07575B}
\definecolor{textgray}{HTML}{34393C}
\definecolor{muted}{HTML}{747A7D}
\color{textgray}
\setstretch{1.13}
\setlength{\parindent}{0pt}
\setlength{\parskip}{5pt}
\setlist[itemize]{leftmargin=7mm,itemsep=2pt,topsep=4pt,parsep=0pt}
\setlist[enumerate]{leftmargin=8.5mm,itemsep=9pt,topsep=6pt,parsep=0pt}
\newcommand{\sect}[1]{\Needspace{8\baselineskip}\vspace{7pt}{\Large\bfseries\color{darkaccent}#1}\par\vspace{3pt}{\color{accent!70}\rule{\linewidth}{0.5pt}}\par\vspace{5pt}}
\newcommand{\subsect}[1]{\Needspace{6\baselineskip}\vspace{4pt}{\normalsize\bfseries\color{darkaccent}#1}\par\vspace{2pt}}
\newcommand{\key}[1]{\textcolor{darkaccent}{\textbf{#1}}}
\newcommand{\checkline}{\par\vspace{2pt}{\small\color{muted}Проверка: знаки $\rightarrow$ скобки $\rightarrow$ арифметика.}\par}
\pagestyle{fancy}
\fancyhf{}
\fancyhead[L]{\small\color{darkaccent} Квадратные уравнения и преобразования}
\fancyhead[R]{\small\color{darkaccent} Володя Хачатурян}
\fancyfoot[C]{\small\color{muted}\thepage}
\renewcommand{\headrulewidth}{0.35pt}
\setlength{\headheight}{14pt}
\begin{document}
\hypersetup{pageanchor=false}
\begin{titlepage}
\thispagestyle{empty}
\begin{center}
\vspace*{28mm}
{\large\color{darkaccent} МАТЕМАТИКА}\par
\vspace{16mm}
{\Huge Чек-лист}\par
\vspace{15mm}
{\LARGE\color{accent} Квадратные уравнения\\[2mm] и алгебраические преобразования}\par
\vspace{14mm}
{\normalsize Дискриминант \textbullet\ Вынесение множителя \textbullet\ Подобные слагаемые}\par
\vfill
{\large Володя Хачатурян}\par
\vspace{12mm}
{\normalsize 2 ноября 2026 г.}\par
\vspace{17mm}
{\normalsize Лёвин Артём Александрович}\\[2mm]
{\small эксклюзивно для Володи Хачатуряна}\par
\vspace{23mm}
\end{center}
\begin{tikzpicture}[remember picture,overlay]
\draw[accent,line width=0.9pt]
  ([xshift=41mm,yshift=25mm]current page.south west)
  .. controls ([xshift=66mm,yshift=37mm]current page.south west) and
  ([xshift=88mm,yshift=14mm]current page.south west)
  .. ([xshift=106mm,yshift=25mm]current page.south west)
  .. controls ([xshift=132mm,yshift=37mm]current page.south west) and
  ([xshift=150mm,yshift=14mm]current page.south west)
  .. ([xshift=-41mm,yshift=25mm]current page.south east);
\end{tikzpicture}
\end{titlepage}
\hypersetup{pageanchor=true}

\sect{1. Знаки и скобки: контроль каждого действия}
Если $a-b=c$, то $b=a-c$. Например,
\[
2-x=10\quad\Longrightarrow\quad -x=8\quad\Longrightarrow\quad\boxed{x=-8}.
\]
\key{Первое правило:} знак относится к самому числу. Если в уравнении $x^2-x-6=0$, то коэффициент при $x$ равен $-1$, свободный член равен $-6$.

\key{Второе правило:} при подстановке отрицательного числа в формулу используйте скобки:
\[
(-1)^2=1,\qquad -1^2=-1,\qquad (-3)^2=9,\qquad -3^2=-9.
\]
Здесь квадрат в записи $-3^2$ относится только к тройке: $-3^2=-(3^2)$.

\subsect{Чётная и нечётная степени}
\[
(-a)^{2n}=a^{2n},\qquad (-a)^{2n+1}=-a^{2n+1}\quad (a>0,\ n=0,1,2,\ldots).
\]
Поэтому $(-2)^{1000}>0$, а $(-2)^{1001}<0$. Нулевая степень числа $-2$ также даёт $1$.

\sect{2. Полное квадратное уравнение}
Квадратное уравнение имеет вид
\[
\boxed{ax^2+bx+c=0,\qquad a\ne0}.
\]
Здесь $a$ --- коэффициент при $x^2$; $b$ --- при $x$; $c$ --- свободный член. Пропущенный член имеет коэффициент $0$.

\subsect{Дискриминант и количество корней}
\[
\boxed{D=b^2-4ac}.
\]
\begin{center}
\begin{tabularx}{\linewidth}{@{}p{26mm}X@{}}
\toprule
Условие & Корни уравнения (действительные числа)\\
\midrule
$D>0$ & Два различных: $\displaystyle x_{1,2}=\frac{-b\pm\sqrt D}{2a}$.\\
$D=0$ & Один корень: $\displaystyle x=-\frac{b}{2a}$ (двукратный).\\
$D<0$ & Действительных корней нет.\\
\bottomrule
\end{tabularx}
\end{center}

\Needspace{15\baselineskip}
\subsect{Пример 1. Два корня}
Решим $x^2-x-6=0$.
\[
\begin{aligned}
a&=1,\qquad b=-1,\qquad c=-6,\\
D&=(-1)^2-4\cdot1\cdot(-6)=1+24=25,\quad \sqrt D=5,\\
x_1&=\frac{-(-1)+5}{2\cdot1}=\frac{6}{2}=3,\\
x_2&=\frac{-(-1)-5}{2\cdot1}=\frac{-4}{2}=-2.
\end{aligned}
\]
\key{Ответ:} $-2$ и $3$. Проверка: $3^2-3-6=0$; $(-2)^2-(-2)-6=0$.

\subsect{Пример 2. Отрицательный старший коэффициент}
Решим $-x^2+3x+4=0$.
\[
\begin{aligned}
a&=-1,\qquad b=3,\qquad c=4,\\
D&=3^2-4\cdot(-1)\cdot4=9+16=25,\\
x_1&=\frac{-3+5}{2\cdot(-1)}=-1,\qquad
x_2=\frac{-3-5}{2\cdot(-1)}=4.
\end{aligned}
\]
\key{Ответ:} $-1$ и $4$. Деление отрицательного числа на отрицательное даёт положительное.

\sect{3. Вынесение общего множителя}
Общий множитель есть во \emph{всех} слагаемых. Его можно вынести за скобки:
\[
ab+ac=a(b+c).
\]
\begin{align*}
5x+5y&=5(x+y),\\
25x+5y&=5(5x+y),\\
5x^2+3x&=x(5x+3),\\
5x^2+25x^3&=5x^2(1+5x).
\end{align*}
При вынесении делите \emph{каждое} слагаемое на общий множитель. Если после деления получилось $1$, в скобках эту единицу нужно сохранить, когда она является самостоятельным слагаемым.

\key{Способ проверки:} раскройте скобки. Вы должны получить исходное выражение.

\sect{4. Уравнение, представленное произведением}
Если произведение равно нулю, то хотя бы один из множителей равен нулю:
\[
AB=0\quad\Longleftrightarrow\quad A=0\ \text{или}\ B=0.
\]
\key{Пример.} Решим $x^2-x=0$:
\[
x(x-1)=0\quad\Longleftrightarrow\quad
\begin{cases}
x=0,\quad\text{или}\\
x-1=0.
\end{cases}
\]
\key{Ответ:} $0$ и $1$.

Точно так же
\[
(x+3)(x+4)(x+10)=0
\quad\Longrightarrow\quad x=-3,\;-4,\;-10.
\]
\key{Важно:} нельзя сразу делить уравнение $x(x-1)=0$ на $x$: так теряется корень $x=0$.

\sect{5. Неполные квадратные уравнения}
У неполного квадратного уравнения отсутствует член с $x$ или свободный член (иногда оба). Часто решение короче, чем через дискриминант.

\subsect{Случай 1: $ax^2+bx=0$}
Вынесите $x$:
\[
ax^2+bx=x(ax+b)=0.
\]
\key{Пример:} $x^2-3x=x(x-3)=0$. \quad\key{Ответ:} $0$ и $3$.

\subsect{Случай 2: $x^2=m$}
\[
\begin{array}{rcll}
m>0 &\Longrightarrow& x=\pm\sqrt m &\text{(два корня)},\\[2pt]
m=0 &\Longrightarrow& x=0 &\text{(один корень)},\\[2pt]
m<0 &\Longrightarrow& \text{корней нет} &\text{(в действительных числах)}.
\end{array}
\]
\key{Примеры:} $x^2-4=0\Rightarrow x=\pm2$;\quad
$x^2-3=0\Rightarrow x=\pm\sqrt3$;\quad
$x^2+5=0$ действительных корней не имеет.

\key{Почему:} квадрат любого действительного числа неотрицателен: $x^2\ge0$.

\sect{6. Подобные слагаемые}
\key{Подобные слагаемые} имеют одинаковую буквенную часть: те же буквы в тех же степенях. Складываются их коэффициенты.
\[
\begin{aligned}
5x+3x&=8x,\\
-5x^2+3x^2&=-2x^2,\\
4k^2+2k^2&=6k^2.
\end{aligned}
\]
Выражения $5x+3y$ и $5x^2+3x$ \emph{не содержат подобных слагаемых}. Второе можно записать как $x(5x+3)$.

\Needspace{14\baselineskip}
\subsect{Пример. Сначала приведите подобные}
Решим $x^2-5+4x^2=7$.
\[
\begin{aligned}
x^2+4x^2&=7+5,\\
5x^2&=12,\\
x^2&=\frac{12}{5},\\
\boxed{x=\pm\sqrt{\frac{12}{5}}}.
\end{aligned}
\]
При делении на $5$ обязательно делятся \emph{обе части} уравнения. При извлечении квадратного корня из положительного числа получаем два противоположных решения.

\sect{7. Типичные ошибки и самопроверка}
\begin{itemize}
\item \key{Знак коэффициента:} в $x^2-x-6=0$ имеем $b=-1$, $c=-6$.
\item \key{Отрицательное число в квадрате:} $(-1)^2=1$, тогда как $-1^2=-1$.
\item \key{Знаки в дроби:} $(-8):(-2)=4$.
\item \key{Общий множитель:} $2k+4t=2(k+2t)$; после вынесения $2$ каждое слагаемое меняется.
\item \key{Нулевой множитель:} в $x(x-3)=0$ есть два решения, включая ноль.
\item \key{Квадрат и знак числа:} у $x^2=-2$ действительных решений нет.
\item \key{Подобные:} $4k+2k^2$ приводить нельзя: степени различаются.
\end{itemize}
\checkline

\clearpage
\thispagestyle{fancy}
\begin{center}
{\Large\color{darkaccent}\bfseries Алгоритм: как выбрать способ решения}\par
\vspace{4mm}
{\small\color{muted}Для квадратного уравнения $ax^2+bx+c=0$, где $a\ne0$}\par
\vspace{9mm}
\end{center}
\begin{center}
\begin{tikzpicture}[
  >=Latex,
  every node/.style={font=\small,align=center},
  start/.style={ellipse,draw=darkaccent,thin,text width=46mm,inner sep=7pt},
  decision/.style={diamond,aspect=2.35,draw=accent,thin,text width=19mm,inner sep=2.5pt},
  action/.style={rectangle,rounded corners=2pt,draw=darkaccent!65,thin,text width=42mm,minimum height=15mm,inner sep=5pt},
  arrow/.style={->,thin,draw=darkaccent}
]
\node[start] (s) at (0,0) {Перенесите все слагаемые влево\\и приведите подобные};
\node[decision] (c0) at (0,-2.4) {$c=0$?};
\node[action] (fact) at (-5.25,-2.4) {Вынесите $x$:\\$x(ax+b)=0$};
\node[decision] (b0) at (0,-5.0) {$b=0$?};
\node[action] (sq) at (5.2,-5.0) {Получите $x^2=-c/a$.\\Сравните $-c/a$ с нулём};
\node[action] (d) at (0,-7.9) {Найдите $D=b^2-4ac$.\\Используйте формулу корней};
\node[start] (finish) at (0,-10.7) {Запишите все корни\\и проверьте подстановкой};
\draw[arrow] (s)--(c0);
\draw[arrow] (c0.west)--node[above] {да} (fact.east);
\draw[arrow] (c0.south)--node[right] {нет} (b0.north);
\draw[arrow] (b0.east)--node[above] {да} (sq.west);
\draw[arrow] (b0.south)--node[right] {нет} (d.north);
\draw[arrow] (d.south)--(finish.north);
\draw[arrow] (fact.south)--++(0,-7.1)-|(finish.west);
\draw[arrow] (sq.south)--++(0,-4.2)-|(finish.east);
\end{tikzpicture}
\end{center}

\clearpage
\sect{Путь Б. Ежедневная практика}
\subsect{Вторник, 3 ноября 2026 г.}
\textit{10 заданий. Проверяйте знаки, скобки и полноту ответа; решения записывайте в тетрадь.}
\begin{enumerate}[itemsep=23pt,topsep=11pt]
\item Решите уравнение: $7-x=15$.
\item Вычислите: $(-4)^2-4^2+(-3)^2$.
\item Приведите подобные: $6a+2a-3a$.
\item Приведите подобные: $3k^2+8k-5k^2+2k$.
\item Вынесите наибольший общий одночлен: $8t+12t^2$.
\item Вынесите наибольший общий одночлен: $6p^2-9p$.
\item Решите: $x^2-5x=0$.
\item Решите: $x^2-16=0$.
\item Решите: $(x-2)(x+6)=0$.
\item Решите через дискриминант: $x^2-3x-10=0$.
\end{enumerate}
\vfill
\begin{flushright}{\small\color{muted}Выполнено: \rule{14mm}{0.45pt} из 10\qquad Дата: \rule{22mm}{0.45pt}}\end{flushright}
\clearpage
\sect{Путь Б. Ежедневная практика}
\subsect{Среда, 4 ноября 2026 г.}
\textit{10 заданий. Решайте по порядку, постепенно переходя к более сложным заданиям.}
\begin{enumerate}[itemsep=23pt,topsep=11pt]
\item Решите: $3-x=-9$.
\item Вычислите: $(-2)^4-2^4+(-3)^3$.
\item Приведите подобные: $9y^2-4y+2y^2+7y$.
\item Вынесите общий одночлен: $15m^2+10m^3$.
\item Вынесите общий одночлен: $12q^3-18q^2$.
\item Решите: $x^2+7x=0$.
\item Решите: $x^2=7$.
\item Решите в действительных числах: $x^2+6=0$.
\item Решите: $(2x-4)(x+5)=0$.
\item Решите через дискриминант: $x^2+2x-15=0$.
\end{enumerate}
\vfill
\begin{flushright}{\small\color{muted}Выполнено: \rule{14mm}{0.45pt} из 10\qquad Дата: \rule{22mm}{0.45pt}}\end{flushright}
\clearpage
\sect{Путь Б. Ежедневная практика}
\subsect{Четверг, 5 ноября 2026 г.}
\textit{10 заданий. Решайте по порядку, постепенно переходя к более сложным заданиям.}
\begin{enumerate}[itemsep=23pt,topsep=11pt]
\item Решите: $8-2x=20$.
\item Вычислите: $(-3)^2-(-2)^3+(-1)^5$.
\item Приведите подобные: $7r^2-3r+5r-12r^2$.
\item Вынесите общий одночлен: $14s^2-21s^3$.
\item Вынесите общий одночлен: $18n^4+12n^3$.
\item Решите: $2x^2-8x=0$.
\item Решите: $3x^2-27=0$.
\item Решите: $(x+1)(x-4)(x+7)=0$.
\item Решите через дискриминант: $-x^2+5x+6=0$.
\item Приведите подобные и решите: $3x^2+2x^2-20=0$.
\end{enumerate}
\vfill
\begin{flushright}{\small\color{muted}Выполнено: \rule{14mm}{0.45pt} из 10\qquad Дата: \rule{22mm}{0.45pt}}\end{flushright}
\clearpage
\sect{Путь Б. Ежедневная практика}
\subsect{Пятница, 6 ноября 2026 г.}
\textit{10 заданий. Решайте по порядку, постепенно переходя к более сложным заданиям.}
\begin{enumerate}[itemsep=23pt,topsep=11pt]
\item Решите: $5-3x=17$.
\item Вычислите: $-(-3)^2+(-3)^2-(-2)^3$.
\item Приведите подобные: $5u^3-2u^2+9u^2-7u^3$.
\item Вынесите общий одночлен: $24p^2-16p^3$.
\item Вынесите общий одночлен: $21b^3+14b^2-7b$.
\item Решите: $4x^2+12x=0$.
\item Решите в действительных числах: $2x^2+18=0$.
\item Решите: $(3x+6)(2x-10)=0$.
\item Решите через дискриминант: $2x^2-5x-3=0$.
\item Приведите подобные и решите: $x^2+3x^2-4x-8=0$.
\end{enumerate}
\vfill
\begin{flushright}{\small\color{muted}Выполнено: \rule{14mm}{0.45pt} из 10\qquad Дата: \rule{22mm}{0.45pt}}\end{flushright}
\clearpage
\sect{Путь Б. Ежедневная практика}
\subsect{Суббота, 7 ноября 2026 г.}
\textit{10 заданий. Решайте по порядку, постепенно переходя к более сложным заданиям.}
\begin{enumerate}[itemsep=23pt,topsep=11pt]
\item Решите: $12-4x=-8$.
\item Вычислите: $(-4)^3-(-2)^4+3^2$.
\item Приведите подобные: $-6a^2+4a-2a^2-9a$.
\item Вынесите общий одночлен: $18y^4-12y^3+6y^2$.
\item Вынесите общий одночлен: $25k^2-15k^3$.
\item Решите: $5x^2-45=0$.
\item Решите: $-2x^2-8x=0$.
\item Решите: $(x+2)(3x-9)(x-7)=0$.
\item Решите через дискриминант: $-2x^2+x+3=0$.
\item Приведите подобные и решите: $2x^2-5+3x^2=40$.
\end{enumerate}
\vfill
\begin{flushright}{\small\color{muted}Выполнено: \rule{14mm}{0.45pt} из 10\qquad Дата: \rule{22mm}{0.45pt}}\end{flushright}
\clearpage
\sect{Путь Б. Ежедневная практика}
\subsect{Воскресенье, 8 ноября 2026 г.}
\textit{10 заданий. Решайте по порядку, постепенно переходя к более сложным заданиям.}
\begin{enumerate}[itemsep=23pt,topsep=11pt]
\item Решите: $4-5x=29$.
\item Вычислите: $(-2)^5+(-3)^2-(-2)^4$.
\item Приведите подобные: $11m^2-3m^3+5m^3-16m^2$.
\item Вынесите общий одночлен: $20q^3-30q^2+10q$.
\item Вынесите общий одночлен: $-18t^3+12t^2$.
\item Решите: $-3x^2+12x=0$.
\item Решите в действительных числах: $5x^2+20=0$.
\item Решите: $(x-1)(2x+8)(3x-6)=0$.
\item Решите через дискриминант: $3x^2+4x-4=0$.
\item Приведите подобные и решите: $x^2-5x+2x^2+4x-6=0$.
\end{enumerate}
\vfill
\begin{flushright}{\small\color{muted}Выполнено: \rule{14mm}{0.45pt} из 10\qquad Дата: \rule{22mm}{0.45pt}}\end{flushright}
\clearpage
\sect{Ответы к ежедневной практике}
\vspace{1mm}
\textit{Сверяйте ответы после выполнения всех 10 заданий за день.}
\vspace{4mm}
\subsect{Вторник, 3 ноября 2026 г.}
\vspace{2mm}
\begin{center}
\renewcommand{\arraystretch}{1.39}
\begin{tabularx}{\linewidth}{@{}>{\centering\arraybackslash}p{13mm}X@{}}
\toprule
№ & Ответ \\
\midrule
1 & $x=-8$ \\
2 & $9$ \\
3 & $5a$ \\
4 & $-2k^2+10k$ \\
5 & $4t(2+3t)$ \\
6 & $3p(2p-3)$ \\
7 & $x=0,\;5$ \\
8 & $x=-4,\;4$ \\
9 & $x=-6,\;2$ \\
10 & $x=-2,\;5$ \\
\bottomrule
\end{tabularx}
\end{center}
\vspace{7mm}
\subsect{Среда, 4 ноября 2026 г.}
\vspace{2mm}
\begin{center}
\renewcommand{\arraystretch}{1.39}
\begin{tabularx}{\linewidth}{@{}>{\centering\arraybackslash}p{13mm}X@{}}
\toprule
№ & Ответ \\
\midrule
1 & $x=12$ \\
2 & $-27$ \\
3 & $11y^2+3y$ \\
4 & $5m^2(3+2m)$ \\
5 & $6q^2(2q-3)$ \\
6 & $x=-7,\;0$ \\
7 & $x=\pm\sqrt7$ \\
8 & Корней нет ($\varnothing$) \\
9 & $x=-5,\;2$ \\
10 & $x=-5,\;3$ \\
\bottomrule
\end{tabularx}
\end{center}
\vspace{7mm}
\clearpage
\sect{Ответы к ежедневной практике}
\vspace{1mm}
\textit{Сверяйте ответы после выполнения всех 10 заданий за день.}
\vspace{4mm}
\subsect{Четверг, 5 ноября 2026 г.}
\vspace{2mm}
\begin{center}
\renewcommand{\arraystretch}{1.39}
\begin{tabularx}{\linewidth}{@{}>{\centering\arraybackslash}p{13mm}X@{}}
\toprule
№ & Ответ \\
\midrule
1 & $x=-6$ \\
2 & $16$ \\
3 & $-5r^2+2r$ \\
4 & $7s^2(2-3s)$ \\
5 & $6n^3(3n+2)$ \\
6 & $x=0,\;4$ \\
7 & $x=-3,\;3$ \\
8 & $x=-7,\;-1,\;4$ \\
9 & $x=-1,\;6$ \\
10 & $x=-2,\;2$ \\
\bottomrule
\end{tabularx}
\end{center}
\vspace{7mm}
\subsect{Пятница, 6 ноября 2026 г.}
\vspace{2mm}
\begin{center}
\renewcommand{\arraystretch}{1.39}
\begin{tabularx}{\linewidth}{@{}>{\centering\arraybackslash}p{13mm}X@{}}
\toprule
№ & Ответ \\
\midrule
1 & $x=-4$ \\
2 & $8$ \\
3 & $-2u^3+7u^2$ \\
4 & $8p^2(3-2p)$ \\
5 & $7b(3b^2+2b-1)$ \\
6 & $x=-3,\;0$ \\
7 & Корней нет ($\varnothing$) \\
8 & $x=-2,\;5$ \\
9 & $x=-\frac12,\;3$ \\
10 & $x=-1,\;2$ \\
\bottomrule
\end{tabularx}
\end{center}
\vspace{7mm}
\clearpage
\sect{Ответы к ежедневной практике}
\vspace{1mm}
\textit{Сверяйте ответы после выполнения всех 10 заданий за день.}
\vspace{4mm}
\subsect{Суббота, 7 ноября 2026 г.}
\vspace{2mm}
\begin{center}
\renewcommand{\arraystretch}{1.39}
\begin{tabularx}{\linewidth}{@{}>{\centering\arraybackslash}p{13mm}X@{}}
\toprule
№ & Ответ \\
\midrule
1 & $x=5$ \\
2 & $-71$ \\
3 & $-8a^2-5a$ \\
4 & $6y^2(3y^2-2y+1)$ \\
5 & $5k^2(5-3k)$ \\
6 & $x=-3,\;3$ \\
7 & $x=-4,\;0$ \\
8 & $x=-2,\;3,\;7$ \\
9 & $x=-1,\;\frac32$ \\
10 & $x=-3,\;3$ \\
\bottomrule
\end{tabularx}
\end{center}
\vspace{7mm}
\subsect{Воскресенье, 8 ноября 2026 г.}
\vspace{2mm}
\begin{center}
\renewcommand{\arraystretch}{1.39}
\begin{tabularx}{\linewidth}{@{}>{\centering\arraybackslash}p{13mm}X@{}}
\toprule
№ & Ответ \\
\midrule
1 & $x=-5$ \\
2 & $-39$ \\
3 & $2m^3-5m^2$ \\
4 & $10q(2q^2-3q+1)$ \\
5 & $6t^2(2-3t)$ \\
6 & $x=0,\;4$ \\
7 & Корней нет ($\varnothing$) \\
8 & $x=-4,\;1,\;2$ \\
9 & $x=-2,\;\frac23$ \\
10 & $x=\dfrac{1\pm\sqrt{73}}6$ \\
\bottomrule
\end{tabularx}
\end{center}
\vspace{7mm}

\sect{Итоговая памятка}
\begin{tabularx}{\linewidth}{@{}p{61mm}X@{}}
\toprule
Что видно в условии & Полезное действие \\
\midrule
$x^2$ и $x$ при свободном члене $0$ & Вынесите $x$ за скобки. \\
$x^2=m$ & Сравните $m$ с нулём. \\
Произведение равно нулю & Приравняйте каждый множитель к нулю. \\
Полное квадратное уравнение & Выпишите $a,b,c$ и найдите $D$. \\
Одинаковая буквенная часть & Сложите коэффициенты. \\
Отрицательное число подставляется в формулу & Используйте скобки и проверьте знаки. \\
\bottomrule
\end{tabularx}
\vfill
\begin{center}\small\color{muted}Лёвин Артём Александрович эксклюзивно для Володи Хачатуряна\end{center}
\end{document}
